\documentclass[11pt, a4paper]{article}

% --- UNIVERSAL PREAMBLE BLOCK ---
\usepackage[a4paper, top=2.5cm, bottom=2.5cm, left=2cm, right=2cm]{geometry}
\usepackage{fontspec}

\usepackage[english, bidi=basic, provide=*]{babel}

\babelprovide[import, onchar=ids fonts]{english}

% Set default/Latin font to Sans Serif in the main (rm) slot
\babelfont{rm}{Noto Sans}

% Packages for math and tables
\usepackage{amsmath}
\usepackage{amssymb}
\usepackage{booktabs}
\usepackage{array}
\usepackage{float}

% Custom commands for units
\newcommand{\uA}{\,\mu\text{A}}
\newcommand{\uV}{\,\mu\text{V}}
\newcommand{\mV}{\,\text{mV}}
\newcommand{\mA}{\,\text{mA}}
\newcommand{\kohm}{\,\text{k}\Omega}
\newcommand{\pF}{\,\text{pF}}
\newcommand{\ns}{\,\text{ns}}
\newcommand{\MHz}{\,\text{MHz}}
\newcommand{\GHz}{\,\text{GHz}}

\title{EE115A Analog Circuits - Homework 7 Solutions}
\author{}
\date{Due date: 12/11/2025}

\begin{document}

\maketitle

\noindent\textbf{Note:} These solutions are provided for your reference. Please ensure you understand the steps and can reproduce the calculations. As per the instructions, Exercise 5 requires Multisim screenshots which you must generate yourself.

\section*{Exercise 1}

\textbf{Given Parameters:}
\begin{itemize}
    \item Technology: $0.18\,\mu\text{m}$
    \item $k'_n = 4k'_p = 400 \uA/\text{V}^2 \implies k'_p = 100 \uA/\text{V}^2$
    \item $V_{tn} = -V_{tp} = 0.4\,\text{V}$
    \item $V_{DD} = 0.9\,\text{V}, \, V_{SS} = -0.9\,\text{V}$
    \item $I_{REF} = 200 \uA$
    \item $V_{ov} = 0.2\,\text{V}$ for all transistors.
\end{itemize}

\subsection*{(a) DC Design}

\textbf{Current Analysis:}
\begin{itemize}
    \item $Q_8$ and $Q_5$ form a current mirror. $Q_8$ and $Q_7$ form a current mirror.
    \item $I_{D8} = I_{REF} = 200 \uA$.
    \item Therefore, $I_{D5} = I_{D7} = 200 \uA$.
    \item $Q_1$ and $Q_2$ are the differential pair, splitting $I_{D5}$ equally (since A and B are grounded).
    \[ I_{D1} = I_{D2} = I_{D5} / 2 = 100 \uA \]
    \item $Q_3$ and $Q_4$ form a current mirror load. $I_{D3} = I_{D1} = 100 \uA$.
    \[ I_{D4} = I_{D3} = 100 \uA \quad (\text{Matches } I_{D2}) \]
    \item $Q_6$ is the second stage amplifying transistor. The problem states $Q_6$ conducts $200 \uA$.
    \[ I_{D6} = 200 \uA \]
\end{itemize}

\textbf{W/L Calculations:}
Using the saturation current formula:
\[ I_D = \frac{1}{2} k' \left(\frac{W}{L}\right) V_{ov}^2 \implies \left(\frac{W}{L}\right) = \frac{2 I_D}{k' V_{ov}^2} \]

We calculate the constant factor for NMOS and PMOS:
\begin{itemize}
    \item For NMOS ($Q_1, Q_2, Q_5, Q_7, Q_8$):
    \[ \frac{2}{k'_n V_{ov}^2} = \frac{2}{400 \times (0.2)^2} = \frac{2}{16} = 0.125 \, (\uA)^{-1} \]
    \item For PMOS ($Q_3, Q_4, Q_6$):
    \[ \frac{2}{k'_p V_{ov}^2} = \frac{2}{100 \times (0.2)^2} = \frac{2}{4} = 0.5 \, (\uA)^{-1} \]
\end{itemize}

\begin{table}[H]
\centering
\caption{Transistor Sizing Results for Exercise 1}
\begin{tabular}{lccccc}
\toprule
\textbf{Device} & \textbf{Type} & \textbf{$I_D$ ($\mu$A)} & \textbf{$k'$ ($\mu$A/V$^2$)} & \textbf{Calc ($2 I_D / k' V_{ov}^2$)} & \textbf{$W/L$} \\
\midrule
$Q_1$ & NMOS & 100 & 400 & $0.125 \times 100$ & \textbf{12.5} \\
$Q_2$ & NMOS & 100 & 400 & $0.125 \times 100$ & \textbf{12.5} \\
$Q_3$ & PMOS & 100 & 100 & $0.5 \times 100$ & \textbf{50} \\
$Q_4$ & PMOS & 100 & 100 & $0.5 \times 100$ & \textbf{50} \\
$Q_5$ & NMOS & 200 & 400 & $0.125 \times 200$ & \textbf{25} \\
$Q_6$ & PMOS & 200 & 100 & $0.5 \times 200$ & \textbf{100} \\
$Q_7$ & NMOS & 200 & 400 & $0.125 \times 200$ & \textbf{25} \\
$Q_8$ & NMOS & 200 & 400 & $0.125 \times 200$ & \textbf{25} \\
\bottomrule
\end{tabular}
\end{table}

\noindent\textbf{DC Output Voltage ($V_O$):} Ideally, if the circuit is perfectly balanced and $V_{DS6}$ and $V_{DS7}$ adjust to satisfy current continuity, the node voltage is typically centered between the rails. Given $V_{DD}=0.9\text{V}$ and $V_{SS}=-0.9\text{V}$, $\mathbf{V_O \approx 0\,\text{V}}$.

\subsection*{(b) Input Common-Mode Range}
The input common-mode range is limited by keeping $Q_1/Q_2$ in saturation (upper limit) and $Q_5$ in saturation (lower limit).

\textbf{Lower Limit ($V_{ICM,min}$):}
$Q_1$ and $Q_5$ must be in saturation. The limiting factor is usually $Q_5$ entering triode region.
\[ V_{in} - V_{GS1} = V_{S1} \implies V_{S1} \ge V_{SS} + V_{ov5} \]
\[ V_{in} \ge V_{SS} + V_{ov5} + V_{GS1} \]
Where $V_{GS1} = V_{tn} + V_{ov1} = 0.4 + 0.2 = 0.6\,\text{V}$.
\[ V_{ICM,min} = -0.9 + 0.2 + 0.6 = \mathbf{-0.1\,\text{V}} \]

\textbf{Upper Limit ($V_{ICM,max}$):}
$Q_1$ must remain in saturation ($V_{DS1} \ge V_{ov1}$).
\[ V_{D1} = V_{DD} - V_{SG3} = V_{DD} - (|V_{tp}| + V_{ov3}) = 0.9 - (0.4 + 0.2) = 0.3\,\text{V} \]
For $Q_1$ saturation: $V_{D1} \ge V_{G1} - V_{tn}$.
\[ 0.3 \ge V_{ICM} - 0.4 \implies V_{ICM} \le 0.7\,\text{V} \]
\[ V_{ICM,max} = \mathbf{0.7\,\text{V}} \]
\textbf{Range:} $[-0.1\,\text{V} \text{ to } 0.7\,\text{V}]$

\subsection*{(c) Allowable Output Voltage Range}
Limited by $Q_6$ and $Q_7$ remaining in saturation.

\textbf{Maximum ($V_{o,max}$):} $Q_6$ stays in saturation: $V_{SD6} \ge V_{ov6}$.
\[ V_{DD} - V_o \ge 0.2 \implies V_o \le 0.9 - 0.2 = \mathbf{0.7\,\text{V}} \]

\textbf{Minimum ($V_{o,min}$):} $Q_7$ stays in saturation: $V_{DS7} \ge V_{ov7}$.
\[ V_o - V_{SS} \ge 0.2 \implies V_o \ge -0.9 + 0.2 = \mathbf{-0.7\,\text{V}} \]
\textbf{Range:} $[-0.7\,\text{V} \text{ to } 0.7\,\text{V}]$

\subsection*{(d) Voltage Gain}
$V_A = 6\,\text{V}$. Gain formula: $A_v = A_1 \times A_2$.

\textbf{Stage 1 parameters ($Q_1, Q_2, Q_3, Q_4$):}
\begin{itemize}
    \item $I_D = 100 \uA$.
    \item $g_{m1} = \frac{2 I_D}{V_{ov}} = \frac{200 \uA}{0.2\,\text{V}} = 1 \mA/\text{V} = 1 \text{ mS}$.
    \item $r_{o2} = \frac{V_A}{I_D} = \frac{6}{100 \uA} = 60 \kohm$.
    \item $r_{o4} = \frac{V_A}{I_D} = 60 \kohm$.
    \item $R_{out1} = r_{o2} || r_{o4} = 30 \kohm$.
    \item $A_1 = g_{m1} R_{out1} = 1 \text{ mS} \times 30 \kohm = 30\,\text{V/V}$.
\end{itemize}

\textbf{Stage 2 parameters ($Q_6, Q_7$):}
\begin{itemize}
    \item $I_D = 200 \uA$.
    \item $g_{m6} = \frac{2 I_D}{V_{ov}} = \frac{400 \uA}{0.2\,\text{V}} = 2 \mA/\text{V} = 2 \text{ mS}$.
    \item $r_{o6} = \frac{V_A}{I_D} = \frac{6}{200 \uA} = 30 \kohm$.
    \item $r_{o7} = \frac{6}{200 \uA} = 30 \kohm$.
    \item $R_{out2} = r_{o6} || r_{o7} = 15 \kohm$.
    \item $A_2 = -g_{m6} R_{out2} = -2 \text{ mS} \times 15 \kohm = -30\,\text{V/V}$.
\end{itemize}

\textbf{Total Gain:}
\[ A_v = A_1 \times A_2 = 30 \times (-30) = \mathbf{-900\,\text{V/V}} \]
(Magnitude: $59.1\,\text{dB}$)

\section*{Exercise 2}

\textbf{Given Parameters:}
\begin{itemize}
    \item Technology: $65\text{-nm}$
    \item $k'_n = 540 \uA/\text{V}^2$, $k'_p = 100 \uA/\text{V}^2$.
    \item $V_{tn} = -V_{tp} = 0.35\,\text{V}$.
    \item $V_{DD} = 1.2\,\text{V}, V_{SS} = 0\,\text{V}$.
    \item $V_{ov} = 0.15\,\text{V}$.
    \item $I_{D1-4} = 200 \uA$, $I_{D6-7} = 400 \uA$.
\end{itemize}

\subsection*{(a) DC Design}

We determine W/L using $W/L = \frac{2 I_D}{k' V_{ov}^2}$.
$V_{ov}^2 = 0.0225\,\text{V}^2$.

\begin{itemize}
    \item \textbf{For NMOS ($Q_1, Q_2, Q_5, Q_7, Q_8$):}
    Denominator factor $k'_n V_{ov}^2 = 540 \times 0.0225 = 12.15 \uA$.
    For $200 \uA$: $W/L = 2(200)/12.15 = 32.9$.
    For $400 \uA$: $W/L = 2(400)/12.15 = 65.8$.
    \item \textbf{For PMOS ($Q_3, Q_4, Q_6$):}
    Denominator factor $k'_p V_{ov}^2 = 100 \times 0.0225 = 2.25 \uA$.
    For $200 \uA$: $W/L = 2(200)/2.25 = 177.8$.
    For $400 \uA$: $W/L = 2(400)/2.25 = 355.6$.
\end{itemize}

\begin{table}[H]
\centering
\caption{Transistor Sizing Results for Exercise 2}
\begin{tabular}{lccc}
\toprule
\textbf{Transistor} & \textbf{$I_D$ ($\mu$A)} & \textbf{Calculation} & \textbf{$W/L$} \\
\midrule
$Q_1, Q_2$ & 200 & $400 / 12.15$ & \textbf{32.9} \\
$Q_3, Q_4$ & 200 & $400 / 2.25$ & \textbf{177.8} \\
$Q_5$ & 400 & $800 / 12.15$ & \textbf{65.8} \\
$Q_6$ & 400 & $800 / 2.25$ & \textbf{355.6} \\
$Q_7$ & 400 & $800 / 12.15$ & \textbf{65.8} \\
$Q_8$ & 400 (assumed) & $800 / 12.15$ & \textbf{65.8} \\
\bottomrule
\end{tabular}
\end{table}

\subsection*{(b) Input Common-Mode Range}
$V_{GS,n} = 0.35 + 0.15 = 0.5\,\text{V}$. $V_{SG,p} = 0.35 + 0.15 = 0.5\,\text{V}$.

\textbf{Lower Limit:}
\[ V_{ICM,min} = V_{SS} + V_{ov5} + V_{GS1} = 0 + 0.15 + 0.5 = \mathbf{0.65\,\text{V}} \]

\textbf{Upper Limit:}
\[ V_{D1} = V_{DD} - V_{SG3} = 1.2 - 0.5 = 0.7\,\text{V} \]
\[ V_{ICM,max} = V_{D1} + V_{tn} = 0.7 + 0.35 = \mathbf{1.05\,\text{V}} \]
\textbf{Range:} $[0.65\,\text{V} \text{ to } 1.05\,\text{V}]$

\subsection*{(c) Allowable Output Voltage Range}
\textbf{Max:} $V_{DD} - V_{ov6} = 1.2 - 0.15 = \mathbf{1.05\,\text{V}}$.\\
\textbf{Min:} $V_{SS} + V_{ov7} = 0 + 0.15 = \mathbf{0.15\,\text{V}}$.\\
\textbf{Range:} $[0.15\,\text{V} \text{ to } 1.05\,\text{V}]$

\subsection*{(d) Voltage Gain}
$V_A = 2.4\,\text{V}$.

\textbf{Stage 1 ($I_D=200 \uA$):}
$g_{m1} = \frac{2(200 \uA)}{0.15} = 2.67 \mA/\text{V}$.
$r_{o2} = r_{o4} = \frac{2.4}{200 \uA} = 12 \kohm$.
$R_{out1} = 6 \kohm$.
$A_1 = g_{m1} R_{out1} = 2.67\text{m} \times 6\text{k} = 16\,\text{V/V}$.

\textbf{Stage 2 ($I_D=400 \uA$):}
$g_{m6} = \frac{2(400 \uA)}{0.15} = 5.33 \mA/\text{V}$.
$r_{o6} = r_{o7} = \frac{2.4}{400 \uA} = 6 \kohm$.
$R_{out2} = 3 \kohm$.
$A_2 = -g_{m6} R_{out2} = -5.33\text{m} \times 3\text{k} = -16\,\text{V/V}$.

\textbf{Total Gain:} $A_v = 16 \times (-16) = \mathbf{-256\,\text{V/V}}$.

\section*{Exercise 3}

\textbf{Given:}
\begin{itemize}
    \item $0.13\text{-}\mu\text{m}$ technology ($L_{min} = 0.13\,\mu\text{m}$).
    \item $L_{ov} = 0.1 L$.
    \item $\mu_n = 400\,\text{cm}^2/\text{V}\cdot\text{s}$.
    \item $V'_A = 5\,\text{V}/\mu\text{m}$.
    \item $V_{ov} = 0.2\,\text{V}$.
\end{itemize}

\textbf{Formulas:}
\begin{enumerate}
    \item \textbf{Intrinsic Gain ($A_0$):}
    \[ A_0 = g_m r_o = \left(\frac{2 I_D}{V_{ov}}\right) \left(\frac{V_A L}{I_D}\right) = \frac{2 V'_A L}{V_{ov}} \]
    Using $V'_A = 5\,\text{V}/\mu\text{m}$, $V_{ov} = 0.2\,\text{V}$:
    \[ A_0 = \frac{2 \times 5 \times L}{0.2} = 50 L \, (\text{where } L \text{ is in } \mu\text{m}) \]

    \item \textbf{Unity-gain Frequency ($f_T$):}
    \[ f_T \simeq \frac{3 \mu_n V_{ov}}{4 \pi L^2} \]
    Be careful with units. $\mu_n = 400\,\text{cm}^2/\text{Vs} = 400 \times 10^8\,\mu\text{m}^2/\text{Vs}$.
    \[ f_T = \frac{3 \times (400 \times 10^8) \times 0.2}{4 \pi L^2} = \frac{1.91 \times 10^9}{L^2}\,\text{Hz} = \frac{1.91}{L^2}\,\text{GHz} \, (\text{with } L \text{ in } \mu\text{m}) \]
\end{enumerate}

\begin{table}[H]
\centering
\caption{Intrinsic Gain and Unity-Gain Frequency}
\begin{tabular}{lccc}
\toprule
\textbf{Length} & \textbf{L ($\mu$m)} & \textbf{$A_0$ (V/V)} & \textbf{$f_T$ (GHz)} \\
\midrule
$L_{min}$ & 0.13 & $50 \times 0.13 = \mathbf{6.5}$ & $1.91 / 0.13^2 \approx \mathbf{113}$ \\
$2L_{min}$ & 0.26 & $50 \times 0.26 = \mathbf{13}$ & $1.91 / 0.26^2 \approx \mathbf{28.3}$ \\
$3L_{min}$ & 0.39 & $50 \times 0.39 = \mathbf{19.5}$ & $1.91 / 0.39^2 \approx \mathbf{12.6}$ \\
$4L_{min}$ & 0.52 & $50 \times 0.52 = \mathbf{26}$ & $1.91 / 0.52^2 \approx \mathbf{7.1}$ \\
$5L_{min}$ & 0.65 & $50 \times 0.65 = \mathbf{32.5}$ & $1.91 / 0.65^2 \approx \mathbf{4.5}$ \\
\bottomrule
\end{tabular}
\end{table}

\section*{Exercise 4}

\textbf{Given:}
\begin{itemize}
    \item $g_m = 5 \mA/\text{V}$.
    \item $C_{gs} = 5 \pF, C_{gd} = 1 \pF, C_L = 5 \pF$.
    \item $R_{sig} = 10 \kohm, R'_L = 10 \kohm$.
\end{itemize}

\textbf{Open-Circuit Time Constants (OCTC) Method:}

1. \textbf{Resistance seen by $C_{gs}$ ($R_{gs0}$):}
\[ R_{gs0} = R_{sig} = 10 \kohm \]
\[ \tau_{gs} = C_{gs} R_{gs0} = 5 \pF \times 10 \kohm = 50 \ns \]

2. \textbf{Resistance seen by $C_{gd}$ ($R_{gd0}$):}
\[ R_{gd0} = R_{sig} (1 + g_m R'_L) + R'_L \]
\[ g_m R'_L = 5 \text{mS} \times 10 \kohm = 50 \]
\[ R_{gd0} = 10\text{k} (1 + 50) + 10\text{k} = 510\text{k} + 10\text{k} = 520 \kohm \]
\[ \tau_{gd} = C_{gd} R_{gd0} = 1 \pF \times 520 \kohm = 520 \ns \]

3. \textbf{Resistance seen by $C_{L}$ ($R_{L0}$):}
\[ R_{L0} = R'_L = 10 \kohm \]
\[ \tau_{L} = C_{L} R_{L0} = 5 \pF \times 10 \kohm = 50 \ns \]

\textbf{Results:}
\begin{itemize}
    \item \textbf{Total Time Constant:}
    \[ \tau_H = \tau_{gs} + \tau_{gd} + \tau_{L} = 50 + 520 + 50 = \mathbf{620 \ns} \]
    \item \textbf{3-dB Frequency ($f_H$):}
    \[ f_H \approx \frac{1}{2 \pi \tau_H} = \frac{1}{2 \pi (620 \times 10^{-9})} \approx \mathbf{256.7 \text{ kHz}} \]
\end{itemize}

\textbf{Percentage caused by interaction of $C_{gd}$ with input:}
This refers to the Miller effect component in $\tau_{gd}$, specifically the term involving $R_{sig}$.
Term $= C_{gd} R_{sig} (1 + g_m R'_L) = 1 \pF \times 510 \kohm = 510 \ns$.
Percentage $= \frac{510}{620} \approx \mathbf{82.3\%}$.

\textbf{To double $f_H$:}
We need $\tau_H$ to be halved: $\tau_{H,new} = 310 \ns$.
Assume $C$ and $R'_L$ are fixed, we solve for new $R_{sig}$.
\[ \tau_H = C_{gs} R_{sig} + C_{gd} [R_{sig}(1+50) + R'_L] + C_L R'_L \]
\[ 310 \ns = 5\text{p}(R_{sig}) + 1\text{p}(51 R_{sig} + 10\text{k}) + 50\text{n} \]
\[ 310\text{n} = 5\text{p} R_{sig} + 51\text{p} R_{sig} + 10\text{n} + 50\text{n} \]
\[ 250\text{n} = 56\text{p} R_{sig} \]
\[ R_{sig} = \frac{250 \times 10^{-9}}{56 \times 10^{-12}} \approx 4464\,\Omega \approx \mathbf{4.46 \kohm} \]

\section*{Exercise 5}

\textbf{Note:} This section provides theoretical calculations. You must set up the circuit in Multisim as shown in Fig P5 to get the simulation results.

\textbf{Given:}
\begin{itemize}
    \item $g_m = 4 \mA/\text{V}, r_o = 20 \kohm$.
    \item $R_{sig} = 20 \kohm, R_L = 20 \kohm$.
    \item $C_{gs} = 2 \pF, C_{gd} = 0.3 \pF, C_{db} = 0.2 \pF, C_L(total) = 1 \pF$.
\end{itemize}

\subsection*{(a) CS Amplifier}

\textbf{Low-Frequency Gain ($A_M$):}
\[ A_M = \frac{v_o}{v_{sig}} = \left(\frac{R_{in}}{R_{in} + R_{sig}}\right) (-g_m R_{out}) \]
For MOSFET at low frequency, $R_{in} \to \infty$, so $\frac{R_{in}}{R_{in} + R_{sig}} \approx 1$.
\[ R_{out} = r_o || R_L = 20\text{k} || 20\text{k} = 10 \kohm \]
\[ A_M = - (4 \text{ mS})(10 \kohm) = \mathbf{-40\,\text{V/V}} \]

\textbf{Estimate $f_H$ using OCTC:}
\begin{itemize}
    \item $\tau_{gs}$: $R_{gs} = R_{sig} = 20 \kohm$.
    $\tau_{gs} = 2 \pF \times 20 \text{k} = 40 \ns$.
    \item $\tau_{gd}$: $R_{gd} = R_{sig}(1 + g_m R_{out}) + R_{out} = 20\text{k}(1 + 40) + 10\text{k} = 820\text{k} + 10\text{k} = 830 \kohm$.
    $\tau_{gd} = 0.3 \pF \times 830 \text{k} = 249 \ns$.
    \item $\tau_{L}$: $R_{Leq} = R_{out} = 10 \kohm$.
    $\tau_{L} = 1 \pF \times 10 \text{k} = 10 \ns$.
\end{itemize}
\[ \tau_H = 40 + 249 + 10 = 299 \ns \]
\[ f_H \approx \frac{1}{2 \pi (299\text{ns})} \approx \mathbf{532 \text{ kHz}} \]

\textbf{Gain-Bandwidth Product (GBW):}
\[ GBW = |A_M| \times f_H = 40 \times 532 \text{ kHz} = \mathbf{21.28 \MHz} \]

\subsection*{(b) Cascode Amplifier}
A Common-Gate (CG) stage is added. The voltage gain of the first stage becomes small ($A_{v1} \approx -g_m / g_m = -1$), suppressing the Miller effect.

\textbf{New $A_M$:}
The total DC gain of a cascode with $R_L \approx r_o$ is approximately the same as the CS stage because the output resistance is dominated by $R_L$.
$R_{out,cascode} \approx (g_m r_o^2) || R_L$. Since $g_m r_o^2 \approx 1.6 \text{ M}\Omega$ and $R_L = 20 \kohm$, $R_{out} \approx 20 \kohm$.
\[ A_M \approx -g_m R_L = -4 \text{m} \times 20 \text{k} = \mathbf{-80\,\text{V/V}} \]

\textbf{New $f_H$:}
Miller effect is removed.
\begin{itemize}
    \item \textbf{Input Node:} $\tau_{in} = R_{sig} (C_{gs1} + C_{gd1}(1 + \frac{g_{m1}}{g_{m2}})) \approx 20\text{k} (2\text{p} + 0.3\text{p}(2)) = 20\text{k}(2.6\text{p}) = 52 \ns$.
    \item \textbf{Output Node:} This is usually the dominant pole. $R_{out} \approx R_L = 20 \kohm$. Total capacitance at output is $C_L + C_{gd2} = 1 \pF + 0.3 \pF = 1.3 \pF$.
    $\tau_{out} = 20\text{k} \times 1.3\text{p} = 26 \ns$.
    \item \textbf{Intermediate Node:} $R \approx 1/g_{m2} = 250 \Omega$. Capacitance $\approx C_{db1} + C_{gs2} = 0.2\text{p} + 2\text{p} = 2.2\text{p}$. $\tau = 0.55 \ns$ (negligible).
\end{itemize}
\[ \tau_H \approx 52 + 26 = 78 \ns \]
\[ f_H \approx \frac{1}{2 \pi (78\text{ns})} \approx \mathbf{2.04 \MHz} \]

\textbf{New GBW:}
\[ GBW = 80 \times 2.04 \MHz \approx \mathbf{163.2 \MHz} \]

\end{document}